% Fichier engendré par tools/generate-tex.py à partir de 2026.
% Ne pas éditer à la main : tout le texte provient des commentaires du fichier Lean.
\documentclass[11pt,a4paper]{article}

% Compile aussi bien avec pdflatex qu'avec xelatex ou tectonic.
\usepackage{iftex}
\ifPDFTeX
  \usepackage[T1]{fontenc}
  \usepackage[utf8]{inputenc}
\else
  \usepackage{fontspec}
\fi
\usepackage[french]{babel}
\usepackage{amsmath,amssymb,amsthm}
\usepackage[margin=2.5cm]{geometry}
% Les figures : TikZ, pour que leurs étiquettes soient composées dans les
% polices du document. Voir la skill `illustrate-theorem`.
\usepackage{tikz}
\usetikzlibrary{calc}

\usepackage[hidelinks]{hyperref}

\theoremstyle{definition}
\newtheorem{definition}{Définition}
\newtheorem{exemple}{Exemple}
\theoremstyle{plain}
\newtheorem{theoreme}{Théorème}
\newtheorem{lemme}[theoreme]{Lemme}

% Un sujet pose des questions, non des théorèmes : voir exams/README.md.
\theoremstyle{definition}
\newtheorem*{reponse}{Réponse}
%% Ce qui n'est pas démontré : la déclaration Lean porte un `sorry`.
\newtheorem*{enonceadmis}{Énoncé admis}

% Un exemple ne se démontre pas : on explique ce qu'il montre.
\newenvironment{explication}{\begin{proof}[Explication]}{\end{proof}}

% Un énoncé écrit mais non démontré : la preuve Lean est un `sorry`.
\newcommand{\admis}{\par\smallskip\noindent\textit{Admis : l'énoncé est écrit en Lean, sa démonstration reste à faire.}\par}

% Renvoi à la déclaration Lean, une ligne après la démonstration, aligné à droite.
\newcommand{\source}[2]{\par\smallskip\noindent\hfill{\footnotesize\href{#1}{\texttt{#2}}}\par\smallskip}

\title{Sujets 2026, décomposés}
\date{}

\begin{document}
\maketitle
% Les identifiants Lean en machine à écrire ne se coupent pas : on tolère des
% espacements plus lâches plutôt que des lignes qui débordent.
\sloppy

\section{Bac2026MetropoleJour1}

\noindent
Baccalauréat, épreuve de spécialité \textemdash{} voie générale, Métropole, 16 juin 2026 (jour 1).

Chaque question du sujet qui porte une proposition mathématique devient un théorème
Lean, dont la réponse est rédigée ci-dessous en français. Celles qui n'en portent pas \textemdash{} compléter un arbre, lire une pente sur un
graphique, écrire un programme \textemdash{} n'ont pas d'énoncé, et l'index du dossier dit
lesquelles et pourquoi.

Les probabilités sont traitées comme le sujet les traite : sur les nombres. La
construction de l'espace probabilisé n'apporterait rien à des calculs qui portent
sur une loi donnée par un tableau, et le chapitre \og{} Probabilités et statistiques \fg{}
du cours fait ce travail-là pour ses propres énoncés.

\subsection{Exercice 1 \textemdash{} probabilités conditionnelles, espérance, Bienaymé-Tchebychev}

\begin{reponse}
Partie A, question 1a : la probabilité de l'évènement contraire est le
complément à \texttt{1} \textemdash{} si \texttt{75 \%} des familles réservent une cabine, \texttt{25 \%} n'en
réservent pas.
\end{reponse}
\begin{proof}
La probabilité d'un évènement et celle de son contraire ont pour somme $1$ :
\[ 1 - 0{,}75 = 0{,}25 . \]
\end{proof}
\source{https://github.com/Commutator-IO/learn-lean/blob/main/exams/2026/Bac2026MetropoleJour1.lean\#L23}{Bac2026MetropoleJour1.lean\#L23}

\begin{reponse}
Partie A, question 2 : \texttt{P(V \(\cap\) C) = P(V) \(\times\) P\_V(C) = 0,3 \(\times\) 0,8 = 0,24}. C'est la
lecture d'un chemin de l'arbre : les probabilités s'y multiplient.
\end{reponse}
\begin{proof}
Les probabilités se multiplient le long d'un chemin de l'arbre :
\[ P(V \cap C) = P(V) \times P_V(C) = 0{,}3 \times 0{,}8 = 0{,}24 . \]
\end{proof}
\source{https://github.com/Commutator-IO/learn-lean/blob/main/exams/2026/Bac2026MetropoleJour1.lean\#L28}{Bac2026MetropoleJour1.lean\#L28}

\begin{reponse}
Partie A, question 3 : \texttt{P\_C(V) = P(V \(\cap\) C) / P(C) = 0,24 / 0,75 = 0,32}.
\end{reponse}
\begin{proof}
Par définition d'une probabilité conditionnelle :
\[ P_C(V) = \frac{P(V \cap C)}{P(C)} = \frac{0{,}24}{0{,}75} = 0{,}32 . \]
\end{proof}
\source{https://github.com/Commutator-IO/learn-lean/blob/main/exams/2026/Bac2026MetropoleJour1.lean\#L31}{Bac2026MetropoleJour1.lean\#L31}

\begin{reponse}
Partie A, question 4 : \texttt{P\_\(\overline{V}\)(C) = (P(C) − P(V \(\cap\) C)) / P(\(\overline{V}\))}, soit \texttt{0,51/0,7},
c'est-à-dire \texttt{0,73} à \texttt{10\(^{-2}\)} près.

La formule des probabilités totales donne \texttt{P(\(\overline{V}\) \(\cap\) C) = P(C) − P(V \(\cap\) C)} : une famille
qui réserve une cabine réserve ou non un véhicule, et les deux cas s'excluent.
\end{reponse}
\begin{proof}
Les évènements $V$ et $\bar{V}$ forment une partition : une famille qui réserve une cabine
réserve ou non un véhicule, et les deux cas s'excluent. La formule des probabilités
totales donne donc
\[ P(\bar{V} \cap C) = P(C) - P(V \cap C) = 0{,}75 - 0{,}24 = 0{,}51 , \]
puis, en divisant par $P(\bar{V}) = 0{,}7$,
\[ P_{\bar{V}}(C) = \frac{0{,}51}{0{,}7} = 0{,}7285714\ldots \]

L'écart à $0{,}73$ vaut $0{,}00142\ldots$, inférieur à $0{,}005$ : l'arrondi au centième
est bien $0{,}73$.
\end{proof}
\source{https://github.com/Commutator-IO/learn-lean/blob/main/exams/2026/Bac2026MetropoleJour1.lean\#L38}{Bac2026MetropoleJour1.lean\#L38}

\begin{reponse}
Partie B, question 1 : \texttt{E(X) = 96} et \texttt{V(X) = 3 114}.

L'espérance est la moyenne des valeurs pondérée par leurs probabilités ; la variance
s'obtient par König\textendash{}Huygens, \texttt{V(X) = E(X\(^{2}\)) − E(X)\(^{2}\)}, ici \texttt{12 330 − 9 216}.
\end{reponse}
\begin{proof}
L'espérance est la moyenne des valeurs pondérée par leurs probabilités :
\[ E(X) = 0 \times 0{,}19 + 70 \times 0{,}06 + 100 \times 0{,}51 + 170 \times 0{,}24
   = 0 + 4{,}2 + 51 + 40{,}8 = 96 . \]

La variance s'obtient par la formule de König–Huygens, $V(X) = E(X^2) - E(X)^2$. Le moment
d'ordre deux vaut
\[ E(X^2) = 4\,900 \times 0{,}06 + 10\,000 \times 0{,}51 + 28\,900 \times 0{,}24
   = 294 + 5\,100 + 6\,936 = 12\,330 , \]
d'où
\[ V(X) = 12\,330 - 96^2 = 12\,330 - 9\,216 = 3\,114 . \]
\end{proof}
\source{https://github.com/Commutator-IO/learn-lean/blob/main/exams/2026/Bac2026MetropoleJour1.lean\#L46}{Bac2026MetropoleJour1.lean\#L46}

\begin{reponse}
Partie B, question 2b : \texttt{E(Z) = 120} et \texttt{V(Z) = 1 728}.

L'espérance est linéaire, donc se multiplie par \texttt{0,6} ; la variance, quadratique, se
multiplie par \texttt{0,6\(^{2}\) = 0,36}. Les variables \texttt{X} et \texttt{Y} étant indépendantes, leurs
variances s'ajoutent \textemdash{} c'est là que l'hypothèse d'indépendance sert, et nulle part
ailleurs.
\end{reponse}
\begin{proof}
L'espérance est linéaire, donc se multiplie par $0{,}6$ et s'ajoute :
\[ E(Z) = 0{,}6\,\bigl(E(X) + E(Y)\bigr) = 0{,}6 \times (96 + 104) = 0{,}6 \times 200 = 120 . \]

La variance, elle, est quadratique : elle se multiplie par $0{,}6^2 = 0{,}36$. Et les
variables $X$ et $Y$ étant indépendantes, leurs variances s'ajoutent \textemdash{} c'est
là que l'hypothèse d'indépendance sert, et nulle part ailleurs :
\[ V(Z) = 0{,}36 \times (3\,114 + 1\,686) = 0{,}36 \times 4\,800 = 1\,728 . \]
\end{proof}
\source{https://github.com/Commutator-IO/learn-lean/blob/main/exams/2026/Bac2026MetropoleJour1.lean\#L57}{Bac2026MetropoleJour1.lean\#L57}

\begin{reponse}
Partie B, question 3a : la moyenne de \texttt{n} variables indépendantes de même loi a
pour espérance celle de la loi, et pour variance la sienne divisée par \texttt{n}. C'est ce
qui fait qu'une moyenne se resserre quand l'échantillon grandit.
\end{reponse}
\begin{proof}
La somme de $n$ variables de même espérance $120$ a pour espérance $120n$ ; diviser par $n$
donne
\[ E(M_n) = \frac{n \times 120}{n} = 120 . \]

Les variables étant indépendantes, leurs variances s'ajoutent : la somme a pour variance
$1\,728\,n$. Diviser par $n$, c'est multiplier par $\dfrac{1}{n}$, donc multiplier la
variance par $\dfrac{1}{n^2}$ :
\[ V(M_n) = \frac{n \times 1\,728}{n^2} = \frac{1\,728}{n} . \]

Les deux simplifications demandent $n \ne 0$. C'est ce quotient par $n$ qui fait qu'une
moyenne se resserre quand l'échantillon grandit.
\end{proof}
\source{https://github.com/Commutator-IO/learn-lean/blob/main/exams/2026/Bac2026MetropoleJour1.lean\#L64}{Bac2026MetropoleJour1.lean\#L64}

\begin{reponse}
Partie B, question 3b : l'inégalité de Bienaymé-Tchebychev donne
\texttt{P(|M\(_{n}\) − 120| \(\ge\) 6) \(\le\) V(M\(_{n}\))/6\(^{2}\) = 48/n}. Pour que la probabilité complémentaire
atteigne \texttt{0,85}, il faut \texttt{48/n \(\le\) 0,15}, c'est-à-dire \texttt{n \(\ge\) 320}.
\end{reponse}
\begin{proof}
Comme $n$ est un entier strictement positif, $n \times 6^2 > 0$ : on peut multiplier les
deux membres par cette quantité sans changer le sens de l'inégalité, qui devient
\[ 1\,728 \le 0{,}15 \times 36 \times n = 5{,}4\,n . \]

Or $\dfrac{1\,728}{5{,}4} = 320$ exactement, donc cette inégalité équivaut à
$320 \le n$ \textemdash{} et comme $320$ est entier, l'équivalence passe telle quelle des
rationnels aux entiers, dans les deux sens.

Le seuil est ici atteint, non dépassé : $n = 320$ convient.
\end{proof}
\source{https://github.com/Commutator-IO/learn-lean/blob/main/exams/2026/Bac2026MetropoleJour1.lean\#L73}{Bac2026MetropoleJour1.lean\#L73}

\subsection{Exercice 2 \textemdash{} vrai ou faux : espace et dénombrement}

\begin{reponse}
Affirmation 1, \textbf{vraie}. Le plan a pour vecteur normal \texttt{(−1 ; 1 ; −5)}, qui est
exactement le vecteur \texttt{AB} : le plan est donc orthogonal à la droite \texttt{(AB)}. Et le
milieu de \texttt{[AB]}, de coordonnées \texttt{(2,5 ; 0,5 ; −0,5)}, vérifie son équation. C'est
le plan médiateur de \texttt{[AB]}.
\end{reponse}
\begin{proof}
Deux vérifications.

D'abord, le vecteur normal du plan, lu sur les coefficients de son équation, est
$(-1 ; 1 ; -5)$ ; et
\[ \vec{AB} = (2 - 3 ;\ 1 - 0 ;\ -3 - 2) = (-1 ;\ 1 ;\ -5) . \]
Les deux coïncident : le plan est orthogonal à la droite $(AB)$.

Ensuite, le milieu de $[AB]$ a pour coordonnées les demi-sommes,
$\left(\dfrac{5}{2} ; \dfrac{1}{2} ; -\dfrac{1}{2}\right)$, et il vérifie l'équation :
\[ -\frac{5}{2} + \frac{1}{2} - 5 \times \left(-\frac{1}{2}\right) - \frac{1}{2}
   = -\frac{5}{2} + \frac{1}{2} + \frac{5}{2} - \frac{1}{2} = 0 . \]

Orthogonal à $(AB)$ et passant par le milieu de $[AB]$ : c'est le plan médiateur de
$[AB]$, et l'affirmation est vraie.
\end{proof}
\source{https://github.com/Commutator-IO/learn-lean/blob/main/exams/2026/Bac2026MetropoleJour1.lean\#L91}{Bac2026MetropoleJour1.lean\#L91}

\begin{reponse}
Affirmation 2, \textbf{fausse}. Chercher un point commun aux deux droites revient à
résoudre un système de trois équations à deux inconnues. Les première et troisième
donnent \texttt{s = 6/7} ; la deuxième est alors fausse. Les droites ne se coupent pas \textemdash{}
elles ne sont pas non plus parallèles, donc elles sont non coplanaires.
\end{reponse}
\begin{proof}
Raisonnons par l'absurde et supposons qu'il existe des paramètres $t$ et $s$ donnant un
point commun. Les trois coordonnées fournissent trois équations à deux inconnues.

La troisième, $2 - 2t = 2 - 5s$, donne $2t = 5s$, c'est-à-dire $t = 2{,}5\,s$. En la
portant dans la première, $t = 3 - s$, il vient $2{,}5\,s = 3 - s$, donc $s = \dfrac{6}{7}$
et $t = \dfrac{15}{7}$.

La deuxième équation, $-1{,}5 - t = s$, devrait alors être satisfaite ; or
\[ -1{,}5 - \frac{15}{7} = -\frac{51}{14} \ne \frac{6}{7} . \]

Contradiction : les deux droites ne se coupent pas. Leurs vecteurs directeurs n'étant pas
colinéaires, elles ne sont pas non plus parallèles : elles sont non coplanaires, et
l'affirmation est fausse.
\end{proof}
\source{https://github.com/Commutator-IO/learn-lean/blob/main/exams/2026/Bac2026MetropoleJour1.lean\#L100}{Bac2026MetropoleJour1.lean\#L100}

\begin{reponse}
Affirmation 3, \textbf{vraie}. Les vecteurs \texttt{CA} et \texttt{CB} ont la même norme, \texttt{4,5}, et
leur produit scalaire vaut \texttt{6,75} : le cosinus de l'angle en \texttt{C} vaut donc \texttt{1/3}, et
\texttt{arccos(1/3)} vaut \texttt{70,5\(^\circ\)} à un dixième près.

Ce qui se démontre ici est le cosinus ; la conversion en degrés arrondis est une
évaluation numérique, que le sujet demande à la calculatrice.
\end{reponse}
\begin{proof}
Le produit scalaire est la somme des produits des coordonnées :
\[ \vec{CA} \cdot \vec{CB} = 1{,}5 \times 0{,}5 + 3 \times 4 + 3 \times (-2)
   = 0{,}75 + 12 - 6 = 6{,}75 . \]

Les deux vecteurs ont le même carré scalaire, donc la même norme :
\[ \|\vec{CA}\|^2 = 1{,}5^2 + 3^2 + 3^2 = 20{,}25 , \qquad
   \|\vec{CB}\|^2 = 0{,}5^2 + 4^2 + 2^2 = 20{,}25 . \]

Le cosinus de l'angle en $C$ est le quotient du produit scalaire par le produit des
normes, c'est-à-dire ici par ce carré commun :
\[ \cos \widehat{C} = \frac{6{,}75}{20{,}25} = \frac{1}{3} . \]

La conversion $\arccos\left(\tfrac{1}{3}\right) \approx 70{,}5^\circ$ est une évaluation
numérique, que le sujet demande à la calculatrice ; ce qui est démontré ici est le cosinus.
\end{proof}
\source{https://github.com/Commutator-IO/learn-lean/blob/main/exams/2026/Bac2026MetropoleJour1.lean\#L112}{Bac2026MetropoleJour1.lean\#L112}

\begin{reponse}
Affirmation 4, \textbf{vraie}. Le code de la porte \texttt{A} est un arrangement de \texttt{3}
symboles parmi \texttt{8}, soit \texttt{8 \(\times\) 7 \(\times\) 6 = 336} possibilités ; celui de la porte \texttt{B} est
une combinaison de \texttt{4} parmi \texttt{8}, soit \texttt{70}. Titouan a une chance sur \texttt{70}, Clotilde
une sur \texttt{336}.

L'ordre est ce qui sépare les deux comptages, et c'est tout l'objet de la question.
\end{reponse}
\begin{proof}
Le code de la porte $A$ est un arrangement : les symboles sont distincts et l'ordre
compte. Il y a $8$ choix pour le premier, $7$ pour le deuxième, $6$ pour le troisième,
soit
\[ 8 \times 7 \times 6 = 336 . \]

Celui de la porte $B$ est une combinaison : les symboles sont distincts et l'ordre ne
compte pas. Il y en a
\[ \binom{8}{4} = \frac{8 \times 7 \times 6 \times 5}{4 \times 3 \times 2 \times 1} = 70 . \]

Chaque code étant équiprobable, Titouan a une chance sur $70$ et Clotilde une sur $336$ ;
comme $\dfrac{1}{70} > \dfrac{1}{336}$, Titouan a la meilleure chance.

L'ordre est ce qui sépare les deux comptages, et c'est tout l'objet de la question.
\end{proof}
\source{https://github.com/Commutator-IO/learn-lean/blob/main/exams/2026/Bac2026MetropoleJour1.lean\#L125}{Bac2026MetropoleJour1.lean\#L125}

\subsection{Exercice 3 \textemdash{} équation différentielle et suite}

\begin{reponse}
Partie A, question 2 : la fonction \texttt{T(t) = 26 − 8e\textasciicircum{}\{−0,035t\}} est solution de
l'équation différentielle \texttt{y' = −0,035y + 0,91}, et vaut \texttt{18} en \texttt{0}.

En dérivant, \texttt{T'(t) = 0,28 e\textasciicircum{}\{−0,035t\}} ; et
\texttt{−0,035 T(t) + 0,91 = −0,91 + 0,28 e\textasciicircum{}\{−0,035t\} + 0,91}, qui est la même chose.
\end{reponse}
\begin{proof}
La fonction $s \mapsto -0{,}035s$ a pour dérivée $-0{,}035$, donc $s \mapsto
\mathrm{e}^{-0{,}035s}$ a pour dérivée $-0{,}035\,\mathrm{e}^{-0{,}035s}$. En dérivant
terme à terme,
\[ T'(t) = 0 - 8 \times (-0{,}035)\,\mathrm{e}^{-0{,}035t}
   = 0{,}28\,\mathrm{e}^{-0{,}035t} . \]

Il reste à voir que c'est bien le second membre de l'équation :
\[ -0{,}035\,T(t) + 0{,}91
   = -0{,}91 + 0{,}28\,\mathrm{e}^{-0{,}035t} + 0{,}91
   = 0{,}28\,\mathrm{e}^{-0{,}035t} . \]

La condition initiale se lit directement : $T(0) = 26 - 8 = 18$.
\end{proof}
\source{https://github.com/Commutator-IO/learn-lean/blob/main/exams/2026/Bac2026MetropoleJour1.lean\#L136}{Bac2026MetropoleJour1.lean\#L136}

\begin{reponse}
Partie A, question 4 : la température ne dépassera jamais \texttt{28 \(^\circ\)C}, ni même
\texttt{26 \(^\circ\)C} : l'exponentielle est strictement positive, donc \texttt{26 − 8e\textasciicircum{}\{−0,035t\} < 26}.

La valeur \texttt{26} est la limite du modèle, et le système n'y arrive jamais.
\end{reponse}

\begin{figure}[h]
\centering
% Construction : T(t) = 26 − 8e^{−0,035t} sur [0 ; 120]. Abscisse = t/20,
% ordonnée = (T − 16)/2,5, de sorte que 18 \(\mapsto\) 0,8, 26 \(\mapsto\) 4 et 28 \(\mapsto\) 4,8.
% La courbe part de 18 °C et croît vers son asymptote 26 °C, sans jamais
% atteindre le plafond 28 °C.
\begin{tikzpicture}[xscale=1.15, yscale=0.95]
  \draw[->] (0, 0) -- (6.6, 0) node[below] {$t$};
  \draw[->] (0, 0) -- (0, 5.4);
  \draw[dashed] (0, 4) -- (6.3, 4) node[right] {$26$};
  \draw[dotted] (0, 4.8) -- (6.3, 4.8) node[right] {$28$};
  \draw[thick] plot coordinates {(0.000, 0.800) (0.100, 1.016) (0.200, 1.218) (0.300, 1.406) (0.400, 1.581) (0.500, 1.745) (0.600, 1.897) (0.700, 2.040) (0.800, 2.172) (0.900, 2.296) (1.000, 2.411) (1.100, 2.518) (1.200, 2.619) (1.300, 2.712) (1.400, 2.799) (1.500, 2.880) (1.600, 2.956) (1.700, 3.026) (1.800, 3.092) (1.900, 3.154) (2.000, 3.211) (2.100, 3.264) (2.200, 3.314) (2.300, 3.360) (2.400, 3.404) (2.500, 3.444) (2.600, 3.482) (2.700, 3.517) (2.800, 3.549) (2.900, 3.580) (3.000, 3.608) (3.100, 3.635) (3.200, 3.659) (3.300, 3.682) (3.400, 3.704) (3.500, 3.724) (3.600, 3.743) (3.700, 3.760) (3.800, 3.776) (3.900, 3.791) (4.000, 3.805) (4.100, 3.819) (4.200, 3.831) (4.300, 3.842) (4.400, 3.853) (4.500, 3.863) (4.600, 3.872) (4.700, 3.881) (4.800, 3.889) (4.900, 3.896) (5.000, 3.903) (5.100, 3.910) (5.200, 3.916) (5.300, 3.922) (5.400, 3.927) (5.500, 3.932) (5.600, 3.937) (5.700, 3.941) (5.800, 3.945) (5.900, 3.949) (6.000, 3.952)};
  \fill (0, 0.8) circle (1.2pt);
  \node[left] at (0, 0.8) {$18$};
  \node[right] at (4.4, 3.55) {$\mathcal{C}_T$};
\end{tikzpicture}

\caption{La température croît vers son asymptote \(26\,^\circ\mathrm{C}\), toujours en deçà du plafond \(28\,^\circ\mathrm{C}\).}
\end{figure}
\begin{proof}
L'exponentielle est strictement positive, donc $8\,\mathrm{e}^{-0{,}035t} > 0$ et
\[ 26 - 8\,\mathrm{e}^{-0{,}035t} < 26 < 28 . \]

La température ne dépasse donc jamais $28\,^\circ\mathrm{C}$, ni même
$26\,^\circ\mathrm{C}$ : cette dernière valeur est la limite du modèle, et le système n'y
arrive jamais.
\end{proof}
\source{https://github.com/Commutator-IO/learn-lean/blob/main/exams/2026/Bac2026MetropoleJour1.lean\#L151}{Bac2026MetropoleJour1.lean\#L151}

\begin{reponse}
Partie B, question 1 : \texttt{u\(_{1}\) = 19,72}.
\end{reponse}
\begin{proof}
On applique la relation de récurrence au rang $0$, où $u_0 = 20$ et
$\mathrm{e}^{0} = 1$ :
\[ u_1 = 0{,}965 \times 20 + 0{,}35 + 0{,}07 = 19{,}3 + 0{,}35 + 0{,}07 = 19{,}72 . \]
\end{proof}
\source{https://github.com/Commutator-IO/learn-lean/blob/main/exams/2026/Bac2026MetropoleJour1.lean\#L156}{Bac2026MetropoleJour1.lean\#L156}

\begin{reponse}
Partie B, question 2 : par récurrence, \texttt{u\(_{n}\) > 10} pour tout \texttt{n}.

L'hérédité tient à ce que \texttt{0,965 \(\times\) 10 + 0,35 = 10} exactement : la borne \texttt{10} est le
point fixe de la partie affine de la récurrence, et le terme exponentiel, positif, ne
peut que l'écarter vers le haut.
\end{reponse}
\begin{proof}
Par récurrence sur $n$.

\emph{Initialisation.} $u_0 = 20 > 10$.

\emph{Hérédité.} Supposons $u_k > 10$. L'exponentielle étant strictement positive,
\[ u_{k+1} = 0{,}965\,u_k + 0{,}35 + 0{,}07\,\mathrm{e}^{-0{,}1k}
   > 0{,}965 \times 10 + 0{,}35 = 9{,}65 + 0{,}35 = 10 . \]

L'hérédité tient à ce que $0{,}965 \times 10 + 0{,}35 = 10$ exactement : la borne $10$ est
le point fixe de la partie affine de la récurrence, et le terme exponentiel, positif, ne
peut qu'écarter la suite vers le haut.
\end{proof}
\source{https://github.com/Commutator-IO/learn-lean/blob/main/exams/2026/Bac2026MetropoleJour1.lean\#L164}{Bac2026MetropoleJour1.lean\#L164}

\begin{reponse}
Partie B, question 4 : la limite vérifie \texttt{x = 0,965x + 0,35}, donc vaut \texttt{10}.

La suite étant décroissante et minorée par \texttt{10}, elle converge ; sa limite est point
fixe de la relation de récurrence, dont le terme exponentiel tend vers zéro.
\end{reponse}
\begin{proof}
Si la suite converge, sa limite $l$ est point fixe de la partie affine de la récurrence,
le terme exponentiel tendant vers zéro : $l = 0{,}965\,l + 0{,}35$. En regroupant,
$0{,}035\,l = 0{,}35$, donc
\[ l = \frac{0{,}35}{0{,}035} = 10 . \]

Ce raisonnement suppose la convergence, que la décroissance et la minoration par $10$
établissent \textemdash{} le sujet la fait démontrer à la question précédente.
\end{proof}
\source{https://github.com/Commutator-IO/learn-lean/blob/main/exams/2026/Bac2026MetropoleJour1.lean\#L179}{Bac2026MetropoleJour1.lean\#L179}

\subsection{Exercice 4 \textemdash{} une fonction avec un logarithme}

\begin{reponse}
Partie A, question 1 : la courbe passe par \texttt{A(0 ; 1)}, et \texttt{f(0) = a} puisque
\texttt{ln 1 = 0} : donc \texttt{a = 1}.
\end{reponse}
\begin{proof}
La courbe passe par $A(0 ; 1)$, donc $f(0) = 1$. Or $\ln(0 + 1) = \ln 1 = 0$, de sorte que
le second terme s'annule :
\[ f(0) = a + b \times \frac{0}{1} = a . \]
Donc $a = 1$.
\end{proof}
\source{https://github.com/Commutator-IO/learn-lean/blob/main/exams/2026/Bac2026MetropoleJour1.lean\#L185}{Bac2026MetropoleJour1.lean\#L185}

\begin{reponse}
Partie A, question 3a : \texttt{f'(x) = b(1 − ln(x+1))/(x+1)\(^{2}\)}.

C'est la dérivée d'un quotient : \texttt{(u/v)' = (u'v − uv')/v\(^{2}\)} avec \texttt{u = b ln(x+1)} et
\texttt{v = x+1}, ce qui donne \texttt{(b/(x+1) \(\times\) (x+1) − b ln(x+1))/(x+1)\(^{2}\)}.
\end{reponse}
\begin{proof}
Sur $]-1 ; +\infty[$, on a $x + 1 > 0$, ce qui autorise la division et le logarithme.

La fonction $t \mapsto t + 1$ a pour dérivée $1$, donc $t \mapsto \ln(t+1)$ a pour dérivée
$\dfrac{1}{x+1}$, par dérivation d'une composée.

La dérivée d'un quotient, $\left(\dfrac{u}{v}\right)' = \dfrac{u'v - uv'}{v^2}$, avec
$u = b\ln(x+1)$ et $v = x+1$, donne alors
\[ f'(x) = \frac{\dfrac{b}{x+1} \times (x+1) - b\ln(x+1) \times 1}{(x+1)^2}
   = \frac{b\,\bigl(1 - \ln(x+1)\bigr)}{(x+1)^2} , \]
la constante $1$ ajoutée n'ayant pas de dérivée.
\end{proof}
\source{https://github.com/Commutator-IO/learn-lean/blob/main/exams/2026/Bac2026MetropoleJour1.lean\#L192}{Bac2026MetropoleJour1.lean\#L192}

\begin{reponse}
Partie A, question 3b : la valeur \texttt{b = 4} se lit sur la tangente en \texttt{A}, puisque
\texttt{f'(0) = b(1 − ln 1)/1\(^{2}\) = b}. La pente, elle, est lue sur le graphique.
\end{reponse}
\begin{proof}
En $x = 0$, le logarithme vaut $\ln 1 = 0$ et le dénominateur vaut $1$, de sorte que
\[ f'(0) = \frac{b\,(1 - 0)}{1} = b . \]
La pente de la tangente en $A$ étant lue $4$ sur le graphique, on en tire $b = 4$.
\end{proof}
\source{https://github.com/Commutator-IO/learn-lean/blob/main/exams/2026/Bac2026MetropoleJour1.lean\#L210}{Bac2026MetropoleJour1.lean\#L210}

\begin{reponse}
Partie B, question 2 : sur \texttt{]−1 ; +\(\infty\)[}, l'inéquation \texttt{1 − ln(x+1) > 0} équivaut à
\texttt{x < e − 1}. Le logarithme étant strictement croissant, il suffit de composer par
l'exponentielle.
\end{reponse}
\begin{proof}
Sur $]-1 ; +\infty[$, on a $x + 1 > 0$. L'inéquation $1 - \ln(x+1) > 0$ s'écrit
$\ln(x+1) < 1$.

Le logarithme et l'exponentielle sont réciproques l'un de l'autre et strictement
croissants : pour $x + 1 > 0$, l'inégalité $\ln(x+1) < 1$ équivaut à
$x + 1 < \mathrm{e}^1$, c'est-à-dire à
\[ x < \mathrm{e} - 1 . \]

Chacun des deux sens s'obtient en composant par la fonction réciproque, ce qui conserve
l'inégalité stricte.
\end{proof}
\source{https://github.com/Commutator-IO/learn-lean/blob/main/exams/2026/Bac2026MetropoleJour1.lean\#L216}{Bac2026MetropoleJour1.lean\#L216}

\begin{reponse}
Partie B, question 5a : \texttt{\(\int\)\(_{0}\)\(^{2}\) ln(x+1)/(x+1) dx = (ln 3)\(^{2}\)/2}.

L'intégrande est de la forme \texttt{u'u} avec \texttt{u = ln(x+1)} : une primitive est \texttt{u\(^{2}\)/2}, et
il ne reste qu'à évaluer entre les bornes, où \texttt{ln 1 = 0} et \texttt{ln 3} restent.
\end{reponse}
\begin{proof}
L'intégrande est de la forme $u'u$ avec $u(x) = \ln(x+1)$ : une primitive est
$\dfrac{u^2}{2}$.

Vérifions-le. Sur $[0 ; 2]$, on a $x + 1 > 0$, et la dérivée de $t \mapsto \ln(t+1)$ vaut
$\dfrac{1}{x+1}$. La dérivée d'un produit donne alors, pour
$F(t) = \dfrac{\ln(t+1) \times \ln(t+1)}{2}$,
\[ F'(x) = \frac{\dfrac{1}{x+1}\ln(x+1) + \ln(x+1)\,\dfrac{1}{x+1}}{2}
   = \frac{\ln(x+1)}{x+1} . \]

La fonction intégrée est continue sur $[0 ; 2]$, donc intégrable. Le théorème fondamental
de l'analyse donne
\[ \int_0^2 \frac{\ln(x+1)}{x+1}\,\mathrm{d}x = F(2) - F(0)
   = \frac{(\ln 3)^2}{2} - \frac{(\ln 1)^2}{2} = \frac{(\ln 3)^2}{2} , \]
puisque $\ln 1 = 0$.
\end{proof}
\source{https://github.com/Commutator-IO/learn-lean/blob/main/exams/2026/Bac2026MetropoleJour1.lean\#L233}{Bac2026MetropoleJour1.lean\#L233}

\section{Bac2026MetropoleJour2}

\noindent
Baccalauréat, épreuve de spécialité \textemdash{} voie générale, Métropole, 17 juin 2026 (jour 2).

Même principe que le jour 1 : chaque question qui porte une proposition devient un
théorème Lean, et l'index du dossier dit lesquelles n'en portent pas.

La géométrie dans l'espace est traitée en coordonnées, comme le sujet la traite : les
points sont des triplets, l'orthogonalité un produit scalaire nul. C'est fidèle à ce
que l'élève écrit, et suffisant pour tout ce que l'exercice demande.

\subsection{Exercice 1 \textemdash{} géométrie dans l'espace}

\begin{reponse}
Question 2a : le vecteur \texttt{n(1 ; −1 ; 4)} est normal au plan \texttt{(ABC)}, car il est
orthogonal à \texttt{AB(1 ; −3 ; −1)} et à \texttt{AC(−2 ; −2 ; 0)}.

Deux produits scalaires nuls suffisent : \texttt{AB} et \texttt{AC} ne sont pas colinéaires, donc ils
engendrent la direction du plan.
\end{reponse}
\begin{proof}
Le produit scalaire est, dans un repère orthonormé, la somme des produits des coordonnées.
Avec $\vec{n}(1 ; -1 ; 4)$ :
\[ \vec{n} \cdot \vec{AB} = 1 \times 1 + (-1) \times (-3) + 4 \times (-1) = 1 + 3 - 4 = 0 , \]
\[ \vec{n} \cdot \vec{AC} = 1 \times (-2) + (-1) \times (-2) + 4 \times 0 = -2 + 2 + 0 = 0 . \]

Les vecteurs $\vec{AB}$ et $\vec{AC}$ ne sont pas colinéaires, donc ils engendrent la
direction du plan : orthogonal à tous deux, $\vec{n}$ est normal à $(ABC)$.
\end{proof}
\source{https://github.com/Commutator-IO/learn-lean/blob/main/exams/2026/Bac2026MetropoleJour2.lean\#L22}{Bac2026MetropoleJour2.lean\#L22}

\begin{reponse}
Question 2b : l'équation cartésienne du plan est \texttt{x − y + 4z − 5 = 0}.

Les coefficients sont ceux du vecteur normal ; la constante se règle pour que \texttt{A} soit
sur le plan, et l'on vérifie que \texttt{B} et \texttt{C} y sont aussi.
\end{reponse}
\begin{proof}
Un plan de vecteur normal $(1 ; -1 ; 4)$ a une équation de la forme $x - y + 4z + d = 0$ ;
la constante se règle pour que $A$ soit sur le plan, ce qui donne $d = -5$. Les trois
points la vérifient :
\[ 2 - 1 + 4 - 5 = 0 , \qquad 3 + 2 + 0 - 5 = 0 , \qquad 0 + 1 + 4 - 5 = 0 . \]
\end{proof}
\source{https://github.com/Commutator-IO/learn-lean/blob/main/exams/2026/Bac2026MetropoleJour2.lean\#L30}{Bac2026MetropoleJour2.lean\#L30}

\begin{reponse}
Question 4 : le projeté orthogonal de \texttt{D(0 ; 0 ; 2)} sur \texttt{(ABC)} est
\texttt{H(−1/6 ; 1/6 ; 4/3)}.

Deux choses à vérifier, et deux seulement : \texttt{H} appartient au plan, et le vecteur \texttt{DH}
est colinéaire au vecteur normal \textemdash{} ici \texttt{DH = −(1/6) n}.
\end{reponse}
\begin{proof}
Deux choses à vérifier, et deux seulement.

D'abord $H$ appartient au plan :
\[ -\frac{1}{6} - \frac{1}{6} + 4 \times \frac{4}{3} - 5
   = -\frac{1}{3} + \frac{16}{3} - 5 = 5 - 5 = 0 . \]

Ensuite $\vec{DH}$ est colinéaire au vecteur normal :
\[ \vec{DH} = \left(-\frac{1}{6} ;\ \frac{1}{6} ;\ \frac{4}{3} - 2\right)
   = \left(-\frac{1}{6} ;\ \frac{1}{6} ;\ -\frac{2}{3}\right)
   = -\frac{1}{6}\,(1 ;\ -1 ;\ 4) . \]

Un point du plan tel que la droite qui le joint à $D$ soit orthogonale au plan est le
projeté orthogonal de $D$, et il est unique.
\end{proof}
\source{https://github.com/Commutator-IO/learn-lean/blob/main/exams/2026/Bac2026MetropoleJour2.lean\#L40}{Bac2026MetropoleJour2.lean\#L40}

\begin{reponse}
Question 5a : le triangle \texttt{ABC} est isocèle en \texttt{B}, car \texttt{BA} et \texttt{BC} ont le même
carré scalaire, \texttt{11}.
\end{reponse}
\begin{proof}
Un triangle est isocèle en $B$ quand $BA = BC$, ce qui revient à comparer les carrés
scalaires, tous deux positifs :
\[ BA^2 = (2-3)^2 + (1+2)^2 + (1-0)^2 = 1 + 9 + 1 = 11 , \]
\[ BC^2 = (0-3)^2 + (-1+2)^2 + (1-0)^2 = 9 + 1 + 1 = 11 . \]
Ils sont égaux, donc $BA = BC$.
\end{proof}
\source{https://github.com/Commutator-IO/learn-lean/blob/main/exams/2026/Bac2026MetropoleJour2.lean\#L48}{Bac2026MetropoleJour2.lean\#L48}

\begin{reponse}
Question 5b : l'aire du triangle \texttt{ABC} vaut \texttt{3\(\sqrt{2}\)}.

Elle est la moitié de la norme du produit vectoriel \texttt{AB \(\land\) AC = (−2 ; 2 ; −8)}, dont le
carré scalaire vaut \texttt{72} ; et \texttt{\(\sqrt{72}\) = 6\(\sqrt{2}\)}.
\end{reponse}
\begin{proof}
L'aire d'un triangle est la moitié de la norme du produit vectoriel de deux de ses côtés.
Ici $\vec{AB} \wedge \vec{AC} = (-2 ; 2 ; -8)$, de carré scalaire
\[ (-2)^2 + 2^2 + (-8)^2 = 4 + 4 + 64 = 72 . \]

L'aire vaut donc $\dfrac{\sqrt{72}}{2}$, qu'on simplifie en écrivant $72 = 6^2 \times 2$ :
la racine d'un produit de facteurs positifs est le produit des racines, et la racine du
carré d'un nombre positif est ce nombre, donc $\sqrt{72} = 6\sqrt{2}$ et
\[ \frac{\sqrt{72}}{2} = \frac{6\sqrt{2}}{2} = 3\sqrt{2} . \]
\end{proof}
\source{https://github.com/Commutator-IO/learn-lean/blob/main/exams/2026/Bac2026MetropoleJour2.lean\#L57}{Bac2026MetropoleJour2.lean\#L57}

\begin{reponse}
Question 6a : le volume du tétraèdre \texttt{ABCD} vaut \texttt{1}.

La hauteur issue de \texttt{D} est la distance \texttt{DH}, de norme \texttt{(1/6)\(\sqrt{18}\) = \(\sqrt{2}\)/2}, et la base
\texttt{ABC} a pour aire \texttt{3\(\sqrt{2}\)} : le produit \texttt{(1/3) \(\times\) 3\(\sqrt{2}\) \(\times\) (\(\sqrt{2}\)/2)} vaut bien \texttt{1}.
\end{reponse}
\begin{proof}
Le volume d'un tétraèdre est le tiers du produit de l'aire d'une base par la hauteur
correspondante. En prenant pour base $ABC$, d'aire $3\sqrt{2}$, et pour hauteur la distance
$DH = \dfrac{\sqrt{2}}{2}$ :
\[ \frac{1}{3} \times 3\sqrt{2} \times \frac{\sqrt{2}}{2}
   = \frac{\sqrt{2} \times \sqrt{2}}{2} = \frac{2}{2} = 1 , \]
où l'on a utilisé $\sqrt{2} \times \sqrt{2} = 2$.
\end{proof}
\source{https://github.com/Commutator-IO/learn-lean/blob/main/exams/2026/Bac2026MetropoleJour2.lean\#L69}{Bac2026MetropoleJour2.lean\#L69}

\begin{reponse}
Question 6b : l'aire du triangle \texttt{BCD} vaut \texttt{3\(\sqrt{2}\)/2}.

Le volume ne dépend pas de la face choisie comme base : en prenant \texttt{BCD} et la hauteur
issue de \texttt{A}, qui vaut \texttt{\(\sqrt{2}\)}, on retrouve \texttt{1}, ce qui détermine l'aire.
\end{reponse}
\begin{proof}
Le volume ne dépend pas de la face choisie comme base. En prenant $BCD$, d'aire $a$, et la
hauteur issue de $A$, qui vaut $\sqrt{2}$, on retrouve $1$ :
\[ \frac{1}{3}\,a\,\sqrt{2} = 1 . \]

Comme $\sqrt{2} > 0$, on peut diviser :
\[ a = \frac{3}{\sqrt{2}} = \frac{3\sqrt{2}}{\sqrt{2} \times \sqrt{2}}
     = \frac{3\sqrt{2}}{2} , \]
en multipliant haut et bas par $\sqrt{2}$ pour chasser le radical du dénominateur.
\end{proof}
\source{https://github.com/Commutator-IO/learn-lean/blob/main/exams/2026/Bac2026MetropoleJour2.lean\#L77}{Bac2026MetropoleJour2.lean\#L77}

\begin{reponse}
Question 7a : les points \texttt{A}, \texttt{B}, \texttt{C} et \texttt{D\(_{k}\)(0 ; 0 ; k)} sont coplanaires pour
\texttt{k = 5/4} \textemdash{} c'est la valeur qui place \texttt{D\(_{k}\)} sur le plan \texttt{(ABC)}. Le projeté orthogonal
de \texttt{D\(_{k}\)} est alors \texttt{D\(_{k}\)} lui-même.
\end{reponse}
\begin{proof}
Le point $D_k(0 ; 0 ; k)$ est coplanaire avec $A$, $B$ et $C$ exactement quand il vérifie
l'équation du plan $(ABC)$ :
\[ 0 - 0 + 4k - 5 = 0 , \]
d'où $4k = 5$ et $k = \dfrac{5}{4}$.

Le projeté orthogonal de $D_k$ sur le plan est alors $D_k$ lui-même.
\end{proof}
\source{https://github.com/Commutator-IO/learn-lean/blob/main/exams/2026/Bac2026MetropoleJour2.lean\#L87}{Bac2026MetropoleJour2.lean\#L87}

\subsection{Exercice 2 \textemdash{} pollution d'un bassin : modèle discret et modèle continu}

\begin{reponse}
Partie A, question 1 : \texttt{V\(_{1}\) = 6} et \texttt{V\(_{2}\) = 11,97}.
\end{reponse}
\begin{proof}
On applique deux fois la relation de récurrence, en partant de $V_0 = 0$ :
\[ V_1 = 0{,}995 \times 0 + 6 = 6 , \qquad
   V_2 = 0{,}995 \times 6 + 6 = 5{,}97 + 6 = 11{,}97 . \]
\end{proof}
\source{https://github.com/Commutator-IO/learn-lean/blob/main/exams/2026/Bac2026MetropoleJour2.lean\#L92}{Bac2026MetropoleJour2.lean\#L92}

\begin{reponse}
Partie A, question 3 : par récurrence, \texttt{V\(_{n}\) \(\le\) V\(_{n+1}\) \(\le\) 1 200}.

L'encadrement est stable parce que \texttt{1 200} est le point fixe de la récurrence :
\texttt{0,995 \(\times\) 1 200 + 6 = 1 200}. Et l'écart \texttt{V\(_{n+1}\) − V\(_{n}\) = 6 − 0,005 V\(_{n}\)} reste positif tant
que \texttt{V\(_{n}\)} ne dépasse pas \texttt{1 200}.
\end{reponse}
\begin{proof}
On démontre d'abord l'encadrement $0 \le V_n \le 1\,200$, par récurrence.

\emph{Initialisation.} $V_0 = 0$, qui est bien entre $0$ et $1\,200$.

\emph{Hérédité.} Supposons $0 \le V_k \le 1\,200$. Alors $V_{k+1} = 0{,}995\,V_k + 6 \ge 6
\ge 0$ ; et
\[ V_{k+1} \le 0{,}995 \times 1\,200 + 6 = 1\,194 + 6 = 1\,200 , \]
puisque $1\,200$ est précisément le point fixe de la récurrence.

L'encadrement acquis, les trois inégalités demandées suivent. La positivité est immédiate.
La croissance vient de
\[ V_{n+1} - V_n = 6 - 0{,}005\,V_n \ge 6 - 0{,}005 \times 1\,200 = 0 , \]
écart qui reste positif tant que $V_n$ ne dépasse pas $1\,200$. Et la majoration de
$V_{n+1}$ est celle qu'on vient d'établir.
\end{proof}
\source{https://github.com/Commutator-IO/learn-lean/blob/main/exams/2026/Bac2026MetropoleJour2.lean\#L100}{Bac2026MetropoleJour2.lean\#L100}

\begin{reponse}
Partie B, question 1b : la fonction \texttt{v(t) = 1 200(1 − e\textasciicircum{}\{−0,005t\})} est solution de
\texttt{y' = −0,005y + 6}.

En dérivant, \texttt{v'(t) = 6 e\textasciicircum{}\{−0,005t\}} ; et \texttt{−0,005 v(t) + 6 = −6 + 6e\textasciicircum{}\{−0,005t\} + 6}.
\end{reponse}
\begin{proof}
La fonction $s \mapsto -0{,}005s$ a pour dérivée $-0{,}005$, donc $s \mapsto
\mathrm{e}^{-0{,}005s}$ a pour dérivée $-0{,}005\,\mathrm{e}^{-0{,}005s}$. En dérivant,
\[ v'(t) = 1\,200 \times 0{,}005\,\mathrm{e}^{-0{,}005t}
   = 6\,\mathrm{e}^{-0{,}005t} . \]

Et le second membre de l'équation vaut la même chose :
\[ -0{,}005\,v(t) + 6 = -6 + 6\,\mathrm{e}^{-0{,}005t} + 6 = 6\,\mathrm{e}^{-0{,}005t} . \]

La condition initiale se lit directement : $v(0) = 1\,200\,(1 - 1) = 0$.
\end{proof}
\source{https://github.com/Commutator-IO/learn-lean/blob/main/exams/2026/Bac2026MetropoleJour2.lean\#L119}{Bac2026MetropoleJour2.lean\#L119}

\begin{reponse}
Partie B, question 2 : le taux de substance polluante plafonne à
\texttt{1 200 / 30 000 = 4 \%}, en deçà des \texttt{5 \%} qui imposeraient un nettoyage complet.

Le volume ne dépasse jamais \texttt{1 200} litres, puisque l'exponentielle est positive.
\end{reponse}
\begin{proof}
L'exponentielle étant strictement positive, $1 - \mathrm{e}^{-0{,}005t} < 1$ et le volume
ne dépasse jamais $1\,200$ litres. Rapporté aux $30\,000$ litres du bassin, le taux
plafonne donc à
\[ \frac{1\,200}{30\,000} = \frac{4}{100} , \]
soit $4\,\%$ ; et $\dfrac{4}{100} < \dfrac{5}{100}$, donc le seuil qui imposerait un
nettoyage complet n'est jamais atteint.
\end{proof}
\source{https://github.com/Commutator-IO/learn-lean/blob/main/exams/2026/Bac2026MetropoleJour2.lean\#L134}{Bac2026MetropoleJour2.lean\#L134}

\begin{reponse}
Partie B, question 3 : le volume dépasse \texttt{50} litres à partir de
\texttt{t = 200 ln(24/23)} heures.

L'équation \texttt{1 200(1 − e\textasciicircum{}\{−0,005t\}) = 50} donne \texttt{e\textasciicircum{}\{−0,005t\} = 23/24}, d'où le
logarithme.
\end{reponse}
\begin{proof}
De $1\,200\,(1 - \mathrm{e}^{-0{,}005t}) = 50$ on tire $1 - \mathrm{e}^{-0{,}005t} =
\dfrac{50}{1\,200} = \dfrac{1}{24}$, c'est-à-dire
\[ \mathrm{e}^{-0{,}005\,t} = \frac{23}{24} . \]

Le logarithme est réciproque de l'exponentielle : en le prenant des deux côtés,
\[ -0{,}005\,t = \ln\frac{23}{24} = -\ln\frac{24}{23} , \]
la dernière égalité venant de $\ln\left(u^{-1}\right) = -\ln u$. En divisant par
$-0{,}005$,
\[ t = 200\,\ln\frac{24}{23} . \]
\end{proof}
\source{https://github.com/Commutator-IO/learn-lean/blob/main/exams/2026/Bac2026MetropoleJour2.lean\#L142}{Bac2026MetropoleJour2.lean\#L142}

\subsection{Exercice 3 \textemdash{} vrai ou faux : probabilités et dénombrement}

\begin{reponse}
Affirmation 1, \textbf{vraie}. La formule de Bayes donne
\texttt{P\_F(O) = P(O) \(\times\) P\_O(F) / P(F) = 0,52 \(\times\) 0,32 / 0,20 = 0,832}.
\end{reponse}
\begin{proof}
La formule de Bayes renverse le conditionnement :
\[ P_F(O) = \frac{P(O) \times P_O(F)}{P(F)}
   = \frac{0{,}52 \times 0{,}32}{0{,}20} = \frac{0{,}1664}{0{,}20} = 0{,}832 . \]
\end{proof}
\source{https://github.com/Commutator-IO/learn-lean/blob/main/exams/2026/Bac2026MetropoleJour2.lean\#L154}{Bac2026MetropoleJour2.lean\#L154}

\begin{reponse}
Affirmation 2, \textbf{fausse} : la valeur vraie est proche de \texttt{0,96}, non de \texttt{0,4}.

L'espérance vaut \texttt{310} et l'écart type environ \texttt{17} ; le seuil \texttt{340} est donc à près de
deux écarts types au-dessus de la moyenne, ce qui laisse largement plus de \texttt{40 \%} en
dessous.

Énoncé admis. Le vérifier demande d'évaluer une somme de \texttt{341} coefficients binomiaux
de rang \texttt{5 000} : c'est un calcul numérique que ni \texttt{decide} ni \texttt{norm\_num} ne mènent, et
que le sujet fait faire à la calculatrice.
\end{reponse}
\begin{enonceadmis}
La déclaration Lean porte un \texttt{sorry} : rien n'est vérifié ici. Ce qui resterait à
établir est l'encadrement de
\[ P(X \le 340) = \sum_{k=0}^{340} \binom{5\,000}{k}\,0{,}062^{\,k}\,0{,}938^{\,5\,000-k} , \]
somme de $341$ coefficients binomiaux de rang $5\,000$ qu'aucune tactique de calcul ne
mène, et que le sujet fait évaluer à la calculatrice.

L'argument d'ordre de grandeur donné plus haut suffit à répondre à la question posée ; il
ne démontre pas l'inégalité écrite en Lean, et c'est pourquoi elle reste admise.
\end{enonceadmis}
\source{https://github.com/Commutator-IO/learn-lean/blob/main/exams/2026/Bac2026MetropoleJour2.lean\#L165}{Bac2026MetropoleJour2.lean\#L165}

\begin{reponse}
Affirmation 3, \textbf{vraie}. Avec \texttt{E(X) = 310} et
\texttt{V(X) = 5 000 \(\times\) 0,062 \(\times\) 0,938 = 290,78}, l'inégalité de Bienaymé-Tchebychev majore
\texttt{P(|X − 310| \(\ge\) 80)} par \texttt{290,78/6 400 \(\approx\) 0,045}, donc en deçà de \texttt{5 \%}.
\end{reponse}
\begin{proof}
La loi de $X$ étant binomiale de paramètres $5\,000$ et $0{,}062$, son espérance vaut
$np = 310$ et sa variance
\[ V(X) = 5\,000 \times 0{,}062 \times 0{,}938 = 290{,}78 . \]

L'inégalité de Bienaymé-Tchebychev majore alors
\[ P\bigl(|X - 310| \ge 80\bigr) \le \frac{V(X)}{80^2} = \frac{290{,}78}{6\,400}
   = 0{,}04543\ldots , \]
qui est bien inférieur à $\dfrac{5}{100}$.
\end{proof}
\source{https://github.com/Commutator-IO/learn-lean/blob/main/exams/2026/Bac2026MetropoleJour2.lean\#L173}{Bac2026MetropoleJour2.lean\#L173}

\begin{reponse}
Affirmation 4, \textbf{vraie}. Constituer une équipe, c'est choisir \texttt{2} musiciens parmi
\texttt{4} et \texttt{3} non-musiciens parmi \texttt{6}, soit \texttt{6 \(\times\) 20 = 120} équipes.

L'ordre ne compte pas : ce sont bien des combinaisons.
\end{reponse}
\begin{proof}
Constituer une équipe, c'est faire deux choix indépendants : $2$ musiciens parmi $4$, puis
$3$ non-musiciens parmi $6$. L'ordre ne compte pas dans une équipe, ce sont donc des
combinaisons, et les deux comptages se multiplient :
\[ \binom{4}{2} \times \binom{6}{3} = 6 \times 20 = 120 . \]
\end{proof}
\source{https://github.com/Commutator-IO/learn-lean/blob/main/exams/2026/Bac2026MetropoleJour2.lean\#L180}{Bac2026MetropoleJour2.lean\#L180}

\subsection{Exercice 4 \textemdash{} la fonction du logo}

\begin{reponse}
Partie B, question 1a : \texttt{f(x) = (2x − 1)e\textasciicircum{}\{−2x+3\}} s'écrit aussi
\texttt{e\(^{2}\) \(\times\) (2x − 1)/e\textasciicircum{}\{2x−1\}}.

C'est la même chose lue autrement : \texttt{e\textasciicircum{}\{−2x+3\} = e\(^{2}\) \(\times\) e\textasciicircum{}\{−(2x−1)\}}. Cette écriture est
celle qui donne la limite, en faisant apparaître le quotient d'un polynôme par une
exponentielle.
\end{reponse}
\begin{proof}
C'est la même expression lue autrement. La propriété
$\mathrm{e}^{a-b} = \dfrac{\mathrm{e}^a}{\mathrm{e}^b}$, appliquée à
$-2x + 3 = 2 - (2x - 1)$, donne
\[ \mathrm{e}^{-2x+3} = \frac{\mathrm{e}^{2}}{\mathrm{e}^{2x-1}} , \]
d'où
\[ (2x - 1)\,\mathrm{e}^{-2x+3}
   = \mathrm{e}^{2} \times \frac{2x - 1}{\mathrm{e}^{2x-1}} . \]

Cette écriture est celle qui donne la limite en $+\infty$, en faisant apparaître le
quotient d'un polynôme par une exponentielle.
\end{proof}
\source{https://github.com/Commutator-IO/learn-lean/blob/main/exams/2026/Bac2026MetropoleJour2.lean\#L190}{Bac2026MetropoleJour2.lean\#L190}

\begin{reponse}
Partie B, question 2a : \texttt{f'(x) = (−4x + 4)e\textasciicircum{}\{−2x+3\}}.

C'est la dérivée d'un produit : \texttt{2 e\textasciicircum{}\{−2x+3\} + (2x − 1)(−2)e\textasciicircum{}\{−2x+3\}}, dont le facteur
commun se met en évidence. Le signe de \texttt{f'} est donc celui de \texttt{1 − x}.
\end{reponse}
\begin{proof}
La fonction $t \mapsto 2t - 1$ a pour dérivée $2$, et $t \mapsto -2t + 3$ a pour dérivée
$-2$, donc $t \mapsto \mathrm{e}^{-2t+3}$ a pour dérivée $-2\,\mathrm{e}^{-2t+3}$.

La dérivée d'un produit donne alors
\[ f'(x) = 2\,\mathrm{e}^{-2x+3} + (2x - 1) \times (-2)\,\mathrm{e}^{-2x+3}
   = \bigl(2 - 4x + 2\bigr)\mathrm{e}^{-2x+3}
   = (-4x + 4)\,\mathrm{e}^{-2x+3} . \]

L'exponentielle étant strictement positive, le signe de $f'$ est celui de $4 - 4x$,
c'est-à-dire celui de $1 - x$.
\end{proof}
\source{https://github.com/Commutator-IO/learn-lean/blob/main/exams/2026/Bac2026MetropoleJour2.lean\#L204}{Bac2026MetropoleJour2.lean\#L204}

\begin{reponse}
Partie C, question 1 : l'intégrale de \texttt{f} entre \texttt{0,5} et \texttt{3,3} vaut
\texttt{0,5 e\(^{2}\) − 3,3 e\textasciicircum{}\{−3,6\}}, soit \texttt{3,6} au dixième.

L'intégration par parties que demande le sujet revient à reconnaître la primitive
\texttt{F(x) = −x e\textasciicircum{}\{−2x+3\}} : en dérivant, \texttt{−e\textasciicircum{}\{−2x+3\} + 2x e\textasciicircum{}\{−2x+3\} = (2x − 1)e\textasciicircum{}\{−2x+3\}}.
\end{reponse}

\begin{figure}[h]
\centering
% Construction : f(x) = (2x − 1)e^{−2x+3} sur [0,5 ; 3,3]. La fonction s'annule
% en 0,5, atteint son maximum e en x = 1, puis décroît. L'aire hachurée, entre
% ces deux bornes, vaut 0,5e² − 3,3e^{−3,6}, soit 3,6 au dixième.
\begin{tikzpicture}[xscale=1.55, yscale=1.05]
  \draw[->] (0.4, 0) -- (3.6, 0) node[below] {$x$};
  \draw[->] (0.5, -0.15) -- (0.5, 3.15);
  \draw[thin] (0.656, 0) -- (0.656, 1.684);
  \draw[thin] (0.811, 0) -- (0.811, 2.468);
  \draw[thin] (0.967, 0) -- (0.967, 2.712);
  \draw[thin] (1.122, 0) -- (1.122, 2.649);
  \draw[thin] (1.278, 0) -- (1.278, 2.426);
  \draw[thin] (1.433, 0) -- (1.433, 2.133);
  \draw[thin] (1.589, 0) -- (1.589, 1.823);
  \draw[thin] (1.744, 0) -- (1.744, 1.526);
  \draw[thin] (1.900, 0) -- (1.900, 1.258);
  \draw[thin] (2.056, 0) -- (2.056, 1.024);
  \draw[thin] (2.211, 0) -- (2.211, 0.825);
  \draw[thin] (2.367, 0) -- (2.367, 0.660);
  \draw[thin] (2.522, 0) -- (2.522, 0.524);
  \draw[thin] (2.678, 0) -- (2.678, 0.413);
  \draw[thin] (2.833, 0) -- (2.833, 0.324);
  \draw[thin] (2.989, 0) -- (2.989, 0.253);
  \draw[thin] (3.144, 0) -- (3.144, 0.197);
  \draw[thick] plot coordinates {(0.500, 0.000) (0.550, 0.669) (0.600, 1.210) (0.650, 1.642) (0.700, 1.981) (0.750, 2.241) (0.800, 2.433) (0.850, 2.569) (0.900, 2.656) (0.950, 2.704) (1.000, 2.718) (1.050, 2.706) (1.100, 2.671) (1.150, 2.618) (1.200, 2.551) (1.250, 2.473) (1.300, 2.387) (1.350, 2.295) (1.400, 2.199) (1.450, 2.100) (1.500, 2.000) (1.550, 1.900) (1.600, 1.801) (1.650, 1.704) (1.700, 1.609) (1.750, 1.516) (1.800, 1.427) (1.850, 1.341) (1.900, 1.258) (1.950, 1.179) (2.000, 1.104) (2.050, 1.032) (2.100, 0.964) (2.150, 0.899) (2.200, 0.838) (2.250, 0.781) (2.300, 0.727) (2.350, 0.676) (2.400, 0.628) (2.450, 0.583) (2.500, 0.541) (2.550, 0.502) (2.600, 0.465) (2.650, 0.431) (2.700, 0.399) (2.750, 0.369) (2.800, 0.342) (2.850, 0.316) (2.900, 0.292) (2.950, 0.270) (3.000, 0.249) (3.050, 0.230) (3.100, 0.212) (3.150, 0.195) (3.200, 0.180) (3.250, 0.166) (3.300, 0.153)};
  \draw[dotted] (1, 0) -- (1, 2.718);
  \fill (1, 2.718) circle (1.1pt);
  \node[below] at (0.5, -0.08) {$0{,}5$};
  \node[below] at (1, -0.08) {$1$};
  \node[below] at (3.3, -0.08) {$3{,}3$};
  \node[above right] at (1, 2.718) {$\mathrm{e}$};
  \node[right] at (2.4, 0.75) {$\mathcal{C}_f$};
\end{tikzpicture}

\caption{La courbe du logo, et l'aire qu'elle délimite entre \(0{,}5\) et \(3{,}3\).}
\end{figure}
\begin{proof}
L'intégration par parties que demande le sujet revient à reconnaître la primitive
$F(x) = -x\,\mathrm{e}^{-2x+3}$, ce qu'une dérivation vérifie :
\[ F'(x) = -\Bigl(\mathrm{e}^{-2x+3} + x \times (-2)\,\mathrm{e}^{-2x+3}\Bigr)
   = (2x - 1)\,\mathrm{e}^{-2x+3} . \]

La fonction intégrée est continue, donc intégrable sur le segment, et le théorème
fondamental de l'analyse donne
\[ \int_{0{,}5}^{3{,}3} (2x - 1)\,\mathrm{e}^{-2x+3}\,\mathrm{d}x = F(3{,}3) - F(0{,}5) . \]

Or $-2 \times 0{,}5 + 3 = 2$ et $-2 \times 3{,}3 + 3 = -3{,}6$, donc
\[ F(3{,}3) - F(0{,}5) = -3{,}3\,\mathrm{e}^{-3{,}6} + 0{,}5\,\mathrm{e}^{2} . \]

La valeur numérique, $3{,}6$ au dixième, se lit à la calculatrice.
\end{proof}
\source{https://github.com/Commutator-IO/learn-lean/blob/main/exams/2026/Bac2026MetropoleJour2.lean\#L222}{Bac2026MetropoleJour2.lean\#L222}

\section{Brevet2026Metropole}

\noindent
Brevet des collèges \textemdash{} série générale, Métropole, 30 juin 2026.

Chaque question du sujet qui porte une proposition mathématique devient un théorème
Lean, dont la réponse est rédigée ci-dessous en français. Celles qui n'en portent pas \textemdash{} une lecture graphique, une interprétation
à rédiger \textemdash{} n'ont pas d'énoncé, et l'index du dossier dit lesquelles et pourquoi :
c'est une des choses que ce dépôt cherche à mesurer.

Les énoncés sont ceux du sujet, pas leur généralisation. Quand la configuration
géométrique est donnée par une figure, ce qu'elle établit \textemdash{} l'égalité des rapports
de Thalès, le parallélisme \textemdash{} devient une hypothèse : la figure n'est pas
formalisée, l'arithmétique qui en découle l'est.

\subsection{Première partie \textemdash{} automatismes}

\begin{reponse}
Question 1 : \texttt{0,75} s'écrit \texttt{3/4}. Passer d'une écriture décimale à une
fraction, c'est simplifier \texttt{75/100} par \texttt{25}.
\end{reponse}
\begin{proof}
Une écriture décimale se lit comme une fraction de dénominateur une puissance de dix, puis
se simplifie :
\[ 0{,}75 = \frac{75}{100} = \frac{25 \times 3}{25 \times 4} = \frac{3}{4} . \]
\end{proof}
\source{https://github.com/Commutator-IO/learn-lean/blob/main/exams/2026/Brevet2026Metropole.lean\#L23}{Brevet2026Metropole.lean\#L23}

\begin{reponse}
Question 2 : \texttt{−4,7 + 3,5 = −1,2}. La somme de deux relatifs de signes
contraires a le signe du plus grand en valeur absolue.
\end{reponse}
\begin{proof}
Les deux nombres sont de signes contraires : on soustrait les distances à zéro et l'on
garde le signe du plus éloigné.
\[ -4{,}7 + 3{,}5 = -(4{,}7 - 3{,}5) = -1{,}2 . \]
\end{proof}
\source{https://github.com/Commutator-IO/learn-lean/blob/main/exams/2026/Brevet2026Metropole.lean\#L27}{Brevet2026Metropole.lean\#L27}

\begin{reponse}
Question 3 : dans un tableau de proportionnalité où \texttt{6} a pour image \texttt{18},
l'image de \texttt{12} est \texttt{36}. L'égalité des produits en croix suffit à la déterminer.
\end{reponse}
\begin{proof}
Dans un tableau de proportionnalité, les produits en croix sont égaux : $6a = 18 \times 12
= 216$, donc
\[ a = \frac{216}{6} = 36 . \]
\end{proof}
\source{https://github.com/Commutator-IO/learn-lean/blob/main/exams/2026/Brevet2026Metropole.lean\#L31}{Brevet2026Metropole.lean\#L31}

\begin{reponse}
Question 4 : dans un sac de \texttt{10} boules rouges, \texttt{4} bleues et \texttt{6} vertes,
toutes équiprobables, la probabilité de tirer une bleue est \texttt{4/20}, soit \texttt{1/5}.
\end{reponse}
\begin{proof}
Le sac contient $10 + 4 + 6 = 20$ boules, toutes de même chance. La probabilité est le
quotient du nombre de cas favorables par le nombre de cas possibles :
\[ \frac{4}{20} = \frac{4 \times 1}{4 \times 5} = \frac{1}{5} . \]
\end{proof}
\source{https://github.com/Commutator-IO/learn-lean/blob/main/exams/2026/Brevet2026Metropole.lean\#L35}{Brevet2026Metropole.lean\#L35}

\begin{reponse}
Question 5 : la solution de \texttt{10x + 16 = −64} est \texttt{−8}.
\end{reponse}
\begin{proof}
De $10x + 16 = -64$ on tire $10x = -80$, puis $x = -8$.
\end{proof}
\source{https://github.com/Commutator-IO/learn-lean/blob/main/exams/2026/Brevet2026Metropole.lean\#L38}{Brevet2026Metropole.lean\#L38}

\begin{reponse}
Question 6 : la notation scientifique de \texttt{0,004 58} est \texttt{4,58 \(\times\) 10\(^{-3}\)}. Elle
est unique : le facteur est compris entre \texttt{1} et \texttt{10}.
\end{reponse}
\begin{proof}
La notation scientifique écrit le nombre comme le produit d'un facteur compris entre $1$
et $10$ par une puissance de dix :
\[ 0{,}00458 = \frac{4{,}58}{1\,000} = 4{,}58 \times 10^{-3} . \]
Le facteur $4{,}58$ est bien entre $1$ et $10$, ce qui rend l'écriture unique.
\end{proof}
\source{https://github.com/Commutator-IO/learn-lean/blob/main/exams/2026/Brevet2026Metropole.lean\#L42}{Brevet2026Metropole.lean\#L42}

\begin{reponse}
Question 8 : le périmètre d'un rectangle de \texttt{10 mm} sur \texttt{5 mm} vaut \texttt{30 mm}.
La réponse est une longueur, non une aire : c'est ce que les deux réponses en
\texttt{mm\(^{2}\)} proposées cherchent à faire confondre.
\end{reponse}
\begin{proof}
Le périmètre d'un rectangle est le double de la somme de sa longueur et de sa largeur :
\[ 2 \times (10 + 5) = 2 \times 15 = 30 . \]
Le résultat est une longueur, en millimètres, et non une aire.
\end{proof}
\source{https://github.com/Commutator-IO/learn-lean/blob/main/exams/2026/Brevet2026Metropole.lean\#L48}{Brevet2026Metropole.lean\#L48}

\begin{reponse}
Question 9 : dans un triangle rectangle en \texttt{E} où l'hypoténuse \texttt{DF} mesure
\texttt{10 cm} et le côté adjacent \texttt{DE} mesure \texttt{6 cm}, le cosinus de l'angle en \texttt{D} vaut
\texttt{3/5}.

Le cosinus d'un angle aigu est le quotient du côté adjacent par l'hypoténuse : la
question est de reconnaître lequel des trois côtés est l'hypoténuse, et le triangle
\texttt{6, 8, 10} est là pour qu'on puisse se tromper.
\end{reponse}

\begin{figure}[h]
\centering
% Construction : E = (0, 0), F = (4, 0), D = (0, 3) — une unité vaut 2 cm.
% L'angle droit est en E ; l'hypoténuse est [DF], de longueur 10 cm. Le cosinus
% de l'angle en D est le quotient du côté adjacent [DE] par l'hypoténuse.
\begin{tikzpicture}[scale=0.95]
  \coordinate (E) at (0, 0);
  \coordinate (F) at (4, 0);
  \coordinate (D) at (0, 3);
  \draw (E) -- (F) -- (D) -- cycle;
  \draw (0.3, 0) -- (0.3, 0.3) -- (0, 0.3);
  \draw (0, 2.35) arc (-90:-53:0.65);
  \node[below left]  at (E) {$E$};
  \node[below right] at (F) {$F$};
  \node[above left]  at (D) {$D$};
  \node[left]  at (0, 1.5) {$6$ cm};
  \node[below] at (2, 0)   {$8$ cm};
  \node[above right] at (2, 1.5) {$10$ cm};
\end{tikzpicture}

\caption{Le triangle \(DEF\), rectangle en \(E\) : l'hypoténuse est \([DF]\), et non le plus grand des deux autres côtés.}
\end{figure}
\begin{proof}
Le triangle est rectangle en $E$, donc son hypoténuse est $[DF]$ \textemdash{} le côté
opposé à l'angle droit. Le cosinus de l'angle en $D$ est le quotient du côté adjacent par
l'hypoténuse :
\[ \cos \widehat{D} = \frac{DE}{DF} = \frac{6}{10} = \frac{3}{5} . \]

Toute la question est de reconnaître lequel des trois côtés est l'hypoténuse ; le triangle
$6$, $8$, $10$ est là pour qu'on puisse se tromper.
\end{proof}
\source{https://github.com/Commutator-IO/learn-lean/blob/main/exams/2026/Brevet2026Metropole.lean\#L58}{Brevet2026Metropole.lean\#L58}

\subsection{Exercice 1 \textemdash{} statistiques}

\begin{reponse}
Le nombre total de médailles étant la somme des trois colonnes, celle qui
manque s'obtient par soustraction : les Pays-Bas ont \texttt{27} médailles d'or.
\end{reponse}
\begin{proof}
Le total est la somme des trois colonnes : $\text{or} + 17 + 12 = 56$. En retranchant,
\[ \text{or} = 56 - 29 = 27 . \]
\end{proof}
\source{https://github.com/Commutator-IO/learn-lean/blob/main/exams/2026/Brevet2026Metropole.lean\#L65}{Brevet2026Metropole.lean\#L65}

\begin{reponse}
De même pour l'Australie, qui en a \texttt{18}.
\end{reponse}
\begin{proof}
De même : $\text{or} + 17 + 28 = 63$, donc $\text{or} = 63 - 45 = 18$.
\end{proof}
\source{https://github.com/Commutator-IO/learn-lean/blob/main/exams/2026/Brevet2026Metropole.lean\#L68}{Brevet2026Metropole.lean\#L68}

\begin{reponse}
Plus de \texttt{20 \%} des médailles de la Grande-Bretagne sont en bronze : la
proportion vaut \texttt{31/124}, c'est-à-dire exactement un quart.
\end{reponse}
\begin{proof}
La proportion de médailles de bronze est $\dfrac{31}{124}$. Elle dépasse $20\,\% =
\dfrac{1}{5}$, car les produits en croix donnent
\[ 1 \times 124 = 124 < 155 = 5 \times 31 . \]

Et comme $124 = 4 \times 31$, cette proportion vaut exactement
\[ \frac{31}{124} = \frac{31}{4 \times 31} = \frac{1}{4} , \]
c'est-à-dire un quart.
\end{proof}
\source{https://github.com/Commutator-IO/learn-lean/blob/main/exams/2026/Brevet2026Metropole.lean\#L72}{Brevet2026Metropole.lean\#L72}

\begin{reponse}
La médiane des neuf totaux est \texttt{82} : quatre pays en ont moins, quatre en ont
plus. C'est la définition même de la médiane pour un effectif impair, et elle ne
demande pas de trier la série pour être vérifiée.
\end{reponse}
\begin{proof}
Les neuf totaux sont
\[ 220 ;\ 124 ;\ 105 ;\ 56 ;\ 89 ;\ 71 ;\ 82 ;\ 75 ;\ 63 . \]
On compte ceux qui sont strictement inférieurs à $82$ : $56$, $71$, $75$ et $63$, soit
quatre. Puis ceux qui lui sont strictement supérieurs : $220$, $124$, $105$ et $89$, soit
quatre également.

L'effectif étant impair, la valeur qui laisse autant de pays de part et d'autre est la
médiane : $82$. Ce comptage ne demande pas de trier la série.
\end{proof}
\source{https://github.com/Commutator-IO/learn-lean/blob/main/exams/2026/Brevet2026Metropole.lean\#L78}{Brevet2026Metropole.lean\#L78}

\begin{reponse}
Le Brésil passe de \texttt{20} à \texttt{26} médailles d'argent : une augmentation de
\texttt{30 \%}. Le taux se lit sur l'écart rapporté à la valeur de départ, non à celle
d'arrivée.
\end{reponse}
\begin{proof}
Le taux d'évolution est l'écart rapporté à la valeur de départ :
\[ \frac{26 - 20}{20} = \frac{6}{20} = \frac{30}{100} , \]
soit une augmentation de $30\,\%$. C'est bien la valeur de départ qui sert de référence,
non celle d'arrivée.
\end{proof}
\source{https://github.com/Commutator-IO/learn-lean/blob/main/exams/2026/Brevet2026Metropole.lean\#L86}{Brevet2026Metropole.lean\#L86}

\subsection{Exercice 2 \textemdash{} Thalès, triangles semblables, aire}

\begin{reponse}
Question 1 : le triangle \texttt{ABC} de côtés \texttt{6,4}, \texttt{4,8} et \texttt{8} est rectangle en
\texttt{A}, par la réciproque du théorème de Pythagore \textemdash{} \texttt{6,4\(^{2}\) + 4,8\(^{2}\) = 40,96 + 23,04 = 64}.
\end{reponse}
\begin{proof}
On calcule les deux membres :
\[ AB^2 + AC^2 = 6{,}4^2 + 4{,}8^2 = 40{,}96 + 23{,}04 = 64 \]
et $BC^2 = 8^2 = 64$. Ils sont égaux.

Par la réciproque du théorème de Pythagore, le triangle $ABC$ est donc rectangle, et son
angle droit est en $A$ \textemdash{} le sommet opposé au plus grand côté.
\end{proof}
\source{https://github.com/Commutator-IO/learn-lean/blob/main/exams/2026/Brevet2026Metropole.lean\#L92}{Brevet2026Metropole.lean\#L92}

\begin{reponse}
Question 2 : \texttt{DE = 6 cm} et \texttt{AE = 3,6 cm}.

Les droites \texttt{(BC)} et \texttt{(DE)} étant parallèles et les points alignés, le théorème
de Thalès donne l'égalité des trois rapports ; c'est elle qu'on prend en
hypothèse, la figure n'étant pas formalisée. Le rapport vaut \texttt{4,8/6,4 = 0,75}, et
il ne reste qu'à multiplier.
\end{reponse}

\begin{figure}[h]
\centering
% Construction : A = (0, 0), B = (4, 0), C = (0, 3), D = (3, 0), E = (0 ; 2,25).
% Une unité vaut 1,6 cm : AB = 6,4 ; AC = 4,8 ; AD = 4,8 ; AE = 3,6 ; BC = 8 ; DE = 6.
% Les points A, D, B sont alignés, ainsi que A, E, C ; (DE) et (BC) ont la même
% pente, \(-3/4\), et le rapport de Thalès vaut 0,75.
\begin{tikzpicture}[scale=1.05]
  \coordinate (A) at (0, 0);
  \coordinate (B) at (4, 0);
  \coordinate (C) at (0, 3);
  \coordinate (D) at (3, 0);
  \coordinate (E) at (0, 2.25);
  \draw (A) -- (B) -- (C) -- cycle;
  \draw (D) -- (E);
  \draw (0.25, 0) -- (0.25, 0.25) -- (0, 0.25);
  \foreach \p in {A, B, C, D, E} \fill (\p) circle (1.3pt);
  \node[below left]  at (A) {$A$};
  \node[below right] at (B) {$B$};
  \node[above left]  at (C) {$C$};
  \node[below] at (D) {$D$};
  \node[left]  at (E) {$E$};
  \node[below] at (1.5, -0.32) {$4{,}8$};
  \node[left]  at (-0.32, 1.1) {$3{,}6$};
  \node[above right] at (1.4, 1.15) {$6$};
  \node[above right] at (2.1, 1.6) {$8$};
\end{tikzpicture}

\caption{La configuration de Thalès : \(D\) sur \([AB]\), \(E\) sur \([AC]\), et \((DE)\) parallèle à \((BC)\).}
\end{figure}
\begin{proof}
Les droites $(BC)$ et $(DE)$ étant parallèles et les points alignés, le théorème de
Thalès donne l'égalité des trois rapports ; c'est elle qu'on prend en hypothèse, la figure
n'étant pas formalisée.

Le rapport se calcule sur les deux longueurs connues :
\[ \frac{AD}{AB} = \frac{4{,}8}{6{,}4} = \frac{3}{4} = 0{,}75 . \]

Il ne reste qu'à multiplier. De $\dfrac{AE}{AC} = 0{,}75$ avec $AC = 4{,}8$ on tire
\[ AE = 0{,}75 \times 4{,}8 = 3{,}6 , \]
et de $\dfrac{DE}{BC} = 0{,}75$ avec $BC = 8$ on tire $DE = 0{,}75 \times 8 = 6$.
\end{proof}
\source{https://github.com/Commutator-IO/learn-lean/blob/main/exams/2026/Brevet2026Metropole.lean\#L102}{Brevet2026Metropole.lean\#L102}

\begin{reponse}
Question 5 : l'aire du quadrilatère \texttt{BCDE} vaut \texttt{24 cm\(^{2}\)}.

Il se découpe en deux triangles rectangles en \texttt{A} : \texttt{ABC}, de côtés de l'angle
droit \texttt{6,4} et \texttt{4,8}, et \texttt{ADE}, de côtés \texttt{4,8} et \texttt{3,6}. Leurs aires valent
\texttt{15,36} et \texttt{8,64}.
\end{reponse}
\begin{proof}
Le quadrilatère se découpe en deux triangles rectangles en $A$, dont l'aire est le
demi-produit des côtés de l'angle droit :
\[ \frac{6{,}4 \times 4{,}8}{2} + \frac{4{,}8 \times 3{,}6}{2}
   = 15{,}36 + 8{,}64 = 24 . \]
\end{proof}
\source{https://github.com/Commutator-IO/learn-lean/blob/main/exams/2026/Brevet2026Metropole.lean\#L115}{Brevet2026Metropole.lean\#L115}

\subsection{Exercice 3 \textemdash{} volume d'une boule}

\begin{reponse}
Partie B, question 1 : le volume d'une boule de rayon \texttt{2,5 cm} vaut \texttt{65 cm\(^{3}\)}
à l'unité près.

L'énoncé donne la formule \texttt{V = 4\(\pi\)R\(^{3}\)/3} : le sujet ne demande pas de la démontrer,
et ce dépôt ne la démontre pas non plus \textemdash{} c'est un des énoncés admis du chapitre
\og{} Grandeurs et mesures \fg{}. Ce qui se vérifie ici est le calcul, et l'arrondi :
\texttt{4\(\pi\)(2,5)\(^{3}\)/3} vaut environ \texttt{65,45}, dont l'arrondi à l'unité est bien \texttt{65}.
\end{reponse}
\begin{proof}
La formule $V = \dfrac{4\pi R^3}{3}$ est donnée par l'énoncé : c'est un résultat admis, ici
comme au chapitre \og{} Grandeurs et mesures \fg{}. Ce qui se démontre est le calcul et
l'arrondi.

Avec $R = 2{,}5$, on a $R^3 = 15{,}625$, donc
\[ V = \frac{4 \times 15{,}625}{3}\,\pi = \frac{62{,}5}{3}\,\pi . \]

L'encadrement $3{,}1415 < \pi < 3{,}1416$ donne alors
\[ 65{,}447\ldots < V < 65{,}450 , \]
donc $|V - 65| < 0{,}5$ : l'arrondi à l'unité est bien $65$ centimètres cubes.
\end{proof}
\source{https://github.com/Commutator-IO/learn-lean/blob/main/exams/2026/Brevet2026Metropole.lean\#L129}{Brevet2026Metropole.lean\#L129}

\begin{reponse}
Partie B, question 2 : avec \texttt{1 000 cm\(^{3}\)} de plastique et des boules de \texttt{65 cm\(^{3}\)},
on en fabrique \texttt{15} au maximum. Le nombre cherché est un quotient euclidien, non
un quotient décimal : \texttt{15,38} boules n'a pas de sens.
\end{reponse}
\begin{proof}
Le nombre de boules est un quotient euclidien : chaque boule consomme $65$ centimètres
cubes, et il n'y a pas de boule fractionnaire. Or
\[ 15 \times 65 = 975 \le 1\,000 \quad\text{et}\quad 1\,000 < 1\,040 = 16 \times 65 . \]
Quinze boules tiennent donc dans les $1\,000$ centimètres cubes, seize non.
\end{proof}
\source{https://github.com/Commutator-IO/learn-lean/blob/main/exams/2026/Brevet2026Metropole.lean\#L138}{Brevet2026Metropole.lean\#L138}

\begin{reponse}
Partie B, question 3 : à \texttt{0,9 g/cm\(^{3}\)}, une boule de \texttt{65 cm\(^{3}\)} pèse \texttt{58,5 g}.
La masse volumique est un coefficient de proportionnalité entre volume et masse.
\end{reponse}
\begin{proof}
La masse volumique est un coefficient de proportionnalité entre volume et masse :
\[ 0{,}9 \times 65 = 58{,}5 \text{ grammes} . \]
\end{proof}
\source{https://github.com/Commutator-IO/learn-lean/blob/main/exams/2026/Brevet2026Metropole.lean\#L142}{Brevet2026Metropole.lean\#L142}

\subsection{Exercice 4 \textemdash{} sachets de bonbons}

\begin{reponse}
Question 1 : on ne peut pas constituer \texttt{16} sachets, car \texttt{16} ne divise pas
\texttt{140} \textemdash{} il en resterait des bonbons au caramel.
\end{reponse}
\begin{proof}
On effectue la division euclidienne : $140 = 16 \times 8 + 12$. Le reste vaut $12$, qui
n'est pas nul, donc $16$ ne divise pas $140$.

Il resterait douze bonbons au caramel : les sachets ne pourraient pas être identiques.
\end{proof}
\source{https://github.com/Commutator-IO/learn-lean/blob/main/exams/2026/Brevet2026Metropole.lean\#L148}{Brevet2026Metropole.lean\#L148}

\begin{reponse}
Question 2 : la décomposition en facteurs premiers de \texttt{140} est \texttt{2\(^{2}\) \(\times\) 5 \(\times\) 7}.
\end{reponse}
\begin{proof}
On divise successivement par les nombres premiers :
\[ 140 = 2 \times 70 = 2 \times 2 \times 35 = 2^2 \times 5 \times 7 . \]
\end{proof}
\source{https://github.com/Commutator-IO/learn-lean/blob/main/exams/2026/Brevet2026Metropole.lean\#L151}{Brevet2026Metropole.lean\#L151}

\begin{reponse}
Question 3 : le nombre maximal de sachets est \texttt{28}, et chacun contient alors
\texttt{4} bonbons à la fraise et \texttt{5} au caramel.

Le nombre de sachets doit diviser \texttt{112} et \texttt{140} : c'est un diviseur commun, et le
plus grand possible est leur PGCD. Les décompositions \texttt{112 = 2\(^{4}\) \(\times\) 7} et
\texttt{140 = 2\(^{2}\) \(\times\) 5 \(\times\) 7} donnent \texttt{2\(^{2}\) \(\times\) 7 = 28}.
\end{reponse}
\begin{proof}
Le nombre de sachets doit diviser à la fois $112$ et $140$, puisque les bonbons de chaque
sorte se répartissent également : c'est un diviseur commun, et le plus grand possible est
leur PGCD.

L'algorithme d'Euclide le calcule : $140 = 1 \times 112 + 28$, puis $112 = 4 \times 28 +
0$. Le dernier reste non nul est $28$, donc
\[ \mathrm{PGCD}(112, 140) = 28 . \]

On peut le retrouver sur les décompositions $112 = 2^4 \times 7$ et $140 = 2^2 \times 5
\times 7$, en gardant les facteurs communs avec leur plus petit exposant : $2^2 \times 7 =
28$.

Chaque sachet contient alors $\dfrac{112}{28} = 4$ bonbons à la fraise et
$\dfrac{140}{28} = 5$ bonbons au caramel.
\end{proof}
\source{https://github.com/Commutator-IO/learn-lean/blob/main/exams/2026/Brevet2026Metropole.lean\#L159}{Brevet2026Metropole.lean\#L159}

\begin{reponse}
Et \texttt{28} est bien maximal parmi les diviseurs communs : tout diviseur commun de
\texttt{112} et \texttt{140} divise \texttt{28}, donc ne le dépasse pas.
\end{reponse}
\begin{proof}
Soit $d$ un diviseur commun de $112$ et $140$. Deux résultats sur les entiers concluent.

Le premier : tout diviseur commun de deux entiers divise leur PGCD
(\texttt{Nat.dvd\_gcd}). C'est la propriété qui fait du PGCD un \emph{plus grand} diviseur
commun au sens de la divisibilité, et elle se lit sur l'algorithme d'Euclide : à chaque
étape, $\mathrm{PGCD}(a, b) = \mathrm{PGCD}(b, a \bmod b)$, et $d$ qui divise $a$ et $b$
divise aussi $a \bmod b = a - \lfloor a/b \rfloor \, b$ ; le dernier reste non nul, qui est
le PGCD, est donc lui aussi divisible par $d$.

Le second : un diviseur d'un entier strictement positif ne le dépasse pas
(\texttt{Nat.le\_of\_dvd}). En effet, si $28 = d \times k$ avec $28 > 0$, alors $k \ge 1$,
donc $d \times k \ge d$, c'est-à-dire $d \le 28$.

En composant les deux : $d$ divise $\mathrm{PGCD}(112, 140) = 28$, donc $d \le 28$. Aucun
diviseur commun ne dépasse $28$, qui est donc bien le nombre maximal de sachets.
\end{proof}
\source{https://github.com/Commutator-IO/learn-lean/blob/main/exams/2026/Brevet2026Metropole.lean\#L164}{Brevet2026Metropole.lean\#L164}

\vfill
\noindent\rule{0.35\linewidth}{0.4pt}\par
\noindent{\footnotesize Leonardo de Moura et Sebastian Ullrich,
\emph{The Lean~4 Theorem Prover and Programming Language},
\emph{Automated Deduction — CADE~28}, Springer, 2021, p.~625--635,
\href{https://doi.org/10.1007/978-3-030-79876-5_37}{doi:10.1007/978-3-030-79876-5\_37}.\par}

\end{document}
